% ============================================================================
% Fichier  : partie3/P03_C08_normes_droit.tex
% Titre    : Normes et Droit de l'Investigation
% Auteur   : MINKA MI NGUIDJOÏ Thierry Emmanuel
% Version  : 1.0 -- novembre 2026
% Corrigés : corriges/Corr_C08.tex
% ============================================================================
% Textes cités dans leur version officielle (Normes/loi_2010-012, lus le
% 08/11/2026) : loi n° 2010/012, art. 26 al. 4, 29 al. 1 à 3, 30 al. 1,
% 52, 55, 56, 65 al. 2, 69, 72, 91 al. 3. Normes ISO/IEC paraphrasées.
% À vérifier par l'auteur : statut du Cameroun au regard de la Convention de
% Budapest et de la Convention de Malabo ; dispositions du Code de procédure
% pénale sur l'expertise.
% ============================================================================

\chapter{Normes et Droit de l'Investigation}
\label{chap:normes-droit}
\chaptermark{Normes et droit}

\epigraph{%
    « La compétence technique dit ce que l'on peut faire.\\
    Le droit dit ce que l'on a le droit de faire. »%
}{Principe du cours}

\niveauChapitre{Fondamental}{%
    Lecture des Chapitres~\ref{chap:investigation} et~\ref{chap:preuve}
}

\begin{boxobjectifs}
\begin{itemize}
    \item Situer les normes de l'investigation numérique (famille ISO/IEC
          27037 à 27043, ISO/IEC 27035, ISO/IEC 17025, NIST) et dire ce que
          chacune encadre.
    \item Citer les articles de la loi n°~2010/012 qui organisent l'enquête
          (articles~52 à~58).
    \item Qualifier des faits au regard des infractions de la loi
          n°~2010/012.
    \item Identifier les obligations de l'exploitant d'un système
          d'information qui pèsent sur la préparation d'une enquête.
    \item Distinguer les pouvoirs de l'OPJ, de l'expert requis et de
          l'enquêteur privé.
    \item Décrire la voie de l'entraide judiciaire internationale.
\end{itemize}
\end{boxobjectifs}

\bigskip

% ============================================================================
\section*{Avant de commencer : soixante jours au lieu de dix ans}
% ============================================================================

Bafoussam, 30 mars 2025. Me~ETOUNDI, la juriste de NYANGA CRÉDIT, relit avec
Mme~NDOUMBE la loi n°~2010/012 relative à la cybersécurité et à la
cybercriminalité au Cameroun. Elle cherche les infractions que la plainte
peut viser. Elle trouve aussi un article qu'elle n'attendait pas~: la loi
oblige les exploitants de systèmes d'information à conserver les données de
connexion et de trafic pendant dix ans. Les journaux de la passerelle de
l'établissement n'en gardaient que soixante jours.

\begin{boxalert}
\textbf{Violation typique.} Traiter le droit comme une question qui ne se
pose qu'au moment de porter plainte.

\smallskip
\textbf{Règle~:} le droit encadre l'enquête du premier geste au dernier~:
qui peut recueillir quoi, comment saisir, comment requérir un tiers, et ce
que l'organisation devait avoir préparé. L'enquêteur qui l'ignore produit
des pièces fragiles et expose son client.
\end{boxalert}

\separateur

% ============================================================================
\section{Les Normes de l'Investigation}
\label{sec:c08-normes}
% ============================================================================

Les normes ne disent pas le droit~; elles décrivent l'état de l'art. Un juge
ou un expert adverse qui conteste une méthode demandera si elle suit une
pratique reconnue. Les normes fournissent cette référence.

\begin{tableaucours}{L{2.8cm}L{5.4cm}L{4.0cm}}{%
    \tableauHeader{Texte} &
    \tableauHeader{Objet (paraphrasé)} &
    \tableauHeader{Chapitres du cours}
}
ISO/IEC 27037 & Identification, recueil, acquisition et préservation de la
preuve numérique & \ref{chap:preuve}, \ref{chap:acquisition} \\
ISO/IEC 27041 & Assurance que les méthodes d'examen sont adaptées et
validées & \ref{chap:disque} à \ref{chap:reseau} \\
ISO/IEC 27042 & Analyse et interprétation de la preuve numérique &
\ref{chap:disque} à \ref{chap:chronologie} \\
ISO/IEC 27043 & Principes et processus de l'investigation, de la
préparation à la clôture & \ref{chap:investigation} \\
ISO/IEC 27035 & Gestion des incidents de sécurité de l'information &
\ref{chap:investigation}, \ref{chap:enquete} \\
ISO/IEC 17025 & Compétence des laboratoires d'essais et d'étalonnage, base
de l'accréditation des laboratoires forensiques & \ref{chap:rapport} \\
NIST SP 800-86 & Intégration de la forensique à la réponse aux incidents &
\ref{chap:investigation} \\
NIST SP 800-61 & Gestion des incidents de sécurité & \ref{chap:investigation} \\
\end{tableaucours}
\vspace{-0.8em}
\begin{center}{\small\itshape Tableau~8.1. Les normes de l'investigation numérique}\end{center}

\begin{boxinfo}
\textbf{ISO/IEC 27041 et la validation des outils.} La norme
ISO/IEC~27041~\cite{iso27041_2015} demande que l'enquêteur puisse montrer
qu'une méthode fait ce qu'il prétend. En pratique, on teste un outil sur une
image de référence dont le contenu est connu, et l'on garde la trace du
test. Un outil non validé n'est pas interdit, mais son résultat porte une
confiance moindre, et l'enquêteur doit le dire.
\end{boxinfo}

% ============================================================================
\section{Le Cadre Camerounais de l'Enquête}
\label{sec:c08-enquete}
% ============================================================================

La loi n°~2010/012 du 21~décembre 2010 relative à la cybersécurité et à la
cybercriminalité~\cite{loi2010012} contient un chapitre consacré à la
procédure. Ses articles~52 à~58 organisent l'enquête.

\begin{boxjuris}
\textbf{\loicam{n°~2010/012}, article~52}~\cite{loi2010012}~:
\textit{« (1) En cas d'infraction cybernétique, les Officiers de Police
Judiciaire à compétence générale et les agents habilités de l'Agence,
procèdent aux enquêtes conformément aux dispositions du Code de Procédure
Pénale. [\ldots] (3) Les Officiers de Police Judiciaire et les agents
habilités de l'Agence peuvent, lors des investigations, accéder aux moyens
de transport, à tout local à usage professionnel, à l'exclusion des
domiciles privés, en vue de rechercher, de constater les infractions, de
demander la communication de tous les documents professionnels et en
prendre copie, recueillir, sur convocation ou sur place, les renseignements
et justifications. »}

\medskip
\textbf{Portée} (analyse). L'enquête pénale appartient aux OPJ et aux agents
habilités de l'Agence nationale des technologies de l'information et de la
communication (ANTIC). Leurs pouvoirs d'accès s'arrêtent au domicile privé,
qui relève des règles générales de la perquisition.
\end{boxjuris}

\begin{tableaucours}{cL{9.4cm}}{%
    \tableauHeader{Article} &
    \tableauHeader{Objet}
}
52 & Enquêteurs compétents (OPJ, agents habilités de l'ANTIC) et pouvoirs
d'accès \\
53 & Perquisition portant sur des données, copie en présence des
personnes, scellés (Chapitre~\ref{chap:preuve}) \\
54 & Perquisitions et saisies selon le Code de procédure pénale, en tenant
compte du dépérissement des preuves \\
55 & Réquisition d'une personne qualifiée pour obtenir la version en clair
de données transformées~; remise de la convention de déchiffrement \\
56 & Réquisition à expert, selon les règles de la commission d'expert \\
57 & Commissions rogatoires nationales et internationales \\
58 & Obligations des prestataires de cryptographie envers les enquêteurs \\
\end{tableaucours}
\vspace{-0.8em}
\begin{center}{\small\itshape Tableau~8.2. Les articles de procédure de la loi n°~2010/012}\end{center}

\begin{boxjuris}
\textbf{\loicam{n°~2010/012}, article~55}~\cite{loi2010012}~:
\textit{« (1) Lorsqu'il apparaît que les données saisies ou obtenues au
cours de l'enquête ou de l'instruction ont fait l'objet d'opérations de
transformation empêchant d'accéder en clair ou sont de nature à
compromettre les informations qu'elles contiennent, le Procureur de la
République, le Juge d'Instruction ou la juridiction de jugement peuvent
réquisitionner toute personne physique ou morale qualifiée, en vue
d'effectuer les opérations techniques permettant d'obtenir la version en
clair desdites données. (2) Lorsqu'un moyen de cryptographie a été utilisé,
les autorités judiciaires peuvent exiger la convention secrète de
déchiffrement du cryptogramme. »}

\medskip
\textbf{\loicam{n°~2010/012}, article~56}~\cite{loi2010012}~:
\textit{« La réquisition prévue à l'article 50 ci-dessus peut être faite à
tout expert. Dans ce cas, son exécution est faite conformément aux
dispositions du Code de procédure pénale relatives à la commission
d'expert. »}

\medskip
\textbf{Portée} (analyse). L'article~55 fonde la réquisition qui permettrait
d'ouvrir l'archive chiffrée \texttt{cpt.7z} (Chapitre~\ref{chap:disque}).
L'article~56 fonde la désignation de M.~ABOUBAKAR~: l'expert requis agit
selon les règles de l'expertise du Code de procédure pénale [dispositions à
vérifier par le lecteur sur le texte officiel], et non selon un mandat de
l'établissement.
\end{boxjuris}

% ============================================================================
\section{Qualifier les Faits}
\label{sec:c08-infractions}
% ============================================================================

L'enquête technique établit des faits~; la qualification les rattache à des
infractions. Elle appartient au procureur et au juge, mais l'enquêteur doit
savoir quels faits chaque infraction exige, pour ne pas oublier d'établir
l'un d'eux.

\begin{boxjuris}
\textbf{\loicam{n°~2010/012}, article~65, alinéa~2}~\cite{loi2010012}~:
\textit{« Est puni des peines prévues à l'alinéa 1 ci-dessus, tout accès non
autorisé, à l'ensemble ou à une partie d'un réseau de communications
électroniques ou d'un système d'information ou d'un équipement
terminal. »}

\medskip
\textbf{\loicam{n°~2010/012}, article~72}~\cite{loi2010012}~:
\textit{« Est puni des peines prévues par l'article 66 ci-dessus celui qui,
de quelque manière que ce soit, sans droit, introduit, altère, efface, ou
supprime, afin d'obtenir un bénéfice économique, les données électroniques,
de manière à causer un préjudice patrimonial à autrui. »}
\end{boxjuris}

\begin{boxscenario}{Opération NYANGA. Les faits et les textes}
\begin{tableaucours}{L{4.6cm}L{2.4cm}L{5.0cm}}{%
    \tableauHeader{Fait établi} &
    \tableauHeader{Texte possible} &
    \tableauHeader{Élément que l'enquête doit établir}
}
Session du 2~mars sur la passerelle, par contournement de
l'authentification & Art.~65 al.~2, art.~69 & Accès sans droit~;
violation des mesures de sécurité (art.~69) \\
Copie du dossier \texttt{Clients} & Art.~69 & Accès « afin d'obtenir des
informations ou des données » \\
Installation d'un service, copie de la mémoire de \texttt{lsass.exe} &
Art.~71 & Introduction de données en vue de modifier celles du système \\
23 virements avec la clé volée & Art.~72 & Introduction de données sans
droit, bénéfice économique, préjudice patrimonial \\
\end{tableaucours}

\medskip
\textbf{Analyse} (confiance moyenne, qualification réservée au parquet).
Les faits établis aux Chapitres~\ref{chap:disque} à~\ref{chap:chronologie}
couvrent les éléments matériels de ces infractions. L'identité des auteurs,
\NEK{}, conditionne les poursuites, pas la qualification des faits.
\end{boxscenario}

\begin{boxjuris}
\textbf{\loicam{n°~2010/012}, article~69}~\cite{loi2010012}~:
\textit{« Est puni d'un emprisonnement de cinq (05) à dix (10) ans et d'une
amende de 10.000.000 (dix millions) à 100.000.000 (cent millions) F CFA ou
de l'une de ces deux peines seulement, celui qui accède sans droit, et en
violation des mesures de sécurité, à l'ensemble ou à une partie d'un réseau
de communications électroniques, d'un système d'information ou d'un
équipement terminal, afin d'obtenir des informations ou des données, en
relation avec un système d'information connecté à un autre système
d'information. »}

\medskip
\textbf{Portée} (analyse). L'article~69 exige une violation des mesures de
sécurité. Le contournement de l'authentification de la passerelle en est
une~; l'enquête doit l'établir techniquement, ce que fera le
Chapitre~\ref{chap:artefact-vuln}. Le même article interdit à l'enquêteur
privé d'entrer sans autorisation dans un système, fût-ce pour suivre
l'attaquant~: remonter jusqu'au serveur 203.0.113.47 et s'y connecter
serait un accès sans droit.
\end{boxjuris}

% ============================================================================
\section{Les Obligations de l'Exploitant}
\label{sec:c08-obligations}
% ============================================================================

La loi pèse aussi sur la victime, au titre de ses obligations d'exploitant de
système d'information. Trois dispositions intéressent directement la
préparation d'une enquête.

\begin{boxjuris}
\textbf{\loicam{n°~2010/012}, article~29}~\cite{loi2010012}~:
\textit{« (1) Les exploitants des systèmes d'information ont l'obligation
de conserver les données de connexion et de trafic de leurs systèmes
d'information pendant une période de dix (10) ans. (2) Les exploitants des
systèmes d'information sont tenus d'installer des mécanismes de
surveillance de contrôle d'accès aux données de leurs systèmes
d'information. Les données conservées peuvent être accessibles lors des
investigations judiciaires. (3) Les installations des exploitants des
systèmes d'information peuvent faire l'objet de perquisition ou de saisie
sur ordre d'une autorité judiciaire dans les conditions prévues par les lois
et règlements en vigueur. »}

\medskip
\textbf{\loicam{n°~2010/012}, article~26, alinéa~4}~\cite{loi2010012}~:
\textit{« Les plates-formes des systèmes d'information font l'objet de
protection contre d'éventuels rayonnements et des intrusions qui pourraient
compromettre l'intégrité des données transmises et contre toute autre
attaque externe notamment par un système de détection d'intrusions. »}

\medskip
\textbf{\loicam{n°~2010/012}, article~30, alinéa~1}~\cite{loi2010012}~:
\textit{« Les exploitants des systèmes d'information évaluent, révisent leurs
systèmes de sécurité et introduisent en cas de nécessité les modifications
appropriées dans leurs pratiques, mesures et techniques de sécurité en
fonction de l'évolution des technologies. »}
\end{boxjuris}

\begin{boxlecture}{NYANGA CRÉDIT face à ses obligations d'exploitant}
\begin{tableaucours}{L{6.6cm}cL{3.6cm}}{%
    \tableauHeader{Obligation} &
    \tableauHeader{Cote} &
    \tableauHeader{Pièce}
}
Conserver dix ans les données de connexion et de trafic (art.~29 al.~1) &
\NON{} & Passerelle~: 60~jours \\
Surveiller le contrôle d'accès aux données (art.~29 al.~2) & \OUI{} &
Journaux d'authentification du domaine \\
Disposer d'un système de détection d'intrusions (art.~26 al.~4) & \OUI{} &
Sonde installée en 2024 \\
Réviser les mesures selon l'évolution des technologies (art.~30 al.~1) &
\NON{} & Correctif non appliqué en 89~jours \\
Les mécanismes ont reçu le visa de l'Agence (art.~26 al.~3) & \NEK{} &
Aucun document produit \\
\end{tableaucours}

\medskip
\textbf{Lecture} (analyse, confiance moyenne). Deux \OUI{}, deux \NON{}, un
\NEK{}. Les deux \NON{} pèsent sur la victime~: la durée de conservation a
failli effacer le premier accès, et l'absence de correction a ouvert la
porte. Un avocat de la défense, ou un client lésé par la fuite, pourra s'en
servir. Le rapport d'enquête ne les cache pas~; il les sépare des faits
reprochés aux attaquants.
\end{boxlecture}

% ============================================================================
\section{Trois Statuts, Trois Pouvoirs}
\label{sec:c08-statuts}
% ============================================================================

\begin{tableaucours}{L{2.6cm}L{3.2cm}L{3.2cm}L{3.2cm}}{%
    \tableauHeader{} &
    \tableauHeader{OPJ, agent de l'ANTIC} &
    \tableauHeader{Expert requis} &
    \tableauHeader{Enquêteur privé}
}
Fondement & Art.~52 & Art.~56, Code de procédure pénale & Contrat avec
l'organisation \\
Mandant & Procureur, juge & L'autorité qui le requiert & Le client \\
Accès aux locaux & Locaux professionnels (art.~52 al.~3) & Selon sa
mission & Ceux du client seulement \\
Saisie et scellés & Oui (art.~53) & Non, il reçoit des scellés & Non~;
il conserve sous chaîne de possession \\
Requérir un tiers & Oui, par réquisition & Non & Non~; il demande, sans
contrainte \\
Destinataire du rapport & Le parquet & L'autorité requérante & Le client \\
\end{tableaucours}
\vspace{-0.8em}
\begin{center}{\small\itshape Tableau~8.3. Les pouvoirs selon le statut}\end{center}

L'enquêteur privé n'a aucun pouvoir de contrainte. Il travaille sur les
systèmes que son client possède ou dont il peut autoriser l'examen, et il
remet ses pièces à l'OPJ quand une plainte est déposée. C'est pourquoi la
qualité de ses premières acquisitions compte autant~: elles deviendront les
pièces de l'enquête pénale.

% ============================================================================
\section{L'Entraide Internationale}
\label{sec:c08-entraide}
% ============================================================================

Les adresses 203.0.113.18, .47 et .62 de l'affaire, si elles étaient
réelles, appartiendraient très probablement à des hébergeurs étrangers.
Obtenir leurs données suppose l'entraide judiciaire internationale.

\begin{boxjuris}
\textbf{\loicam{n°~2010/012}, article~91, alinéa~3}~\cite{loi2010012}~:
\textit{« En cas d'urgence, les demandes d'entraide demandées par les
autorités camerounaises ou étrangères peuvent être transmises directement
aux autorités de l'Etat requis pour leur exécution. Le renvoi des pièces
d'exécution aux autorités compétentes de l'Etat requérant est effectué selon
les mêmes modalités. »}

\medskip
\textbf{Portée} (analyse). La voie ordinaire passe par le ministère chargé
des relations extérieures (alinéa~1)~; l'urgence autorise la transmission
directe. Le dépérissement des données des hébergeurs justifie souvent
l'urgence.
\end{boxjuris}

Deux conventions internationales structurent la coopération. La Convention
du Conseil de l'Europe sur la cybercriminalité, dite de
Budapest~\cite{budapest2001}, organise la conservation rapide des données et
l'entraide entre ses parties. La Convention de l'Union africaine sur la
cybersécurité et la protection des données à caractère personnel, dite de
Malabo~\cite{ua2014malabo}, entrée en vigueur en 2023, fixe un cadre
africain. Le statut du Cameroun au regard de chacune [à vérifier par le
lecteur sur les sites des deux organisations] commande la voie utilisable~:
sans convention applicable, l'entraide repose sur la réciprocité et les
accords bilatéraux.

% ============================================================================
\section{Ce que le Droit demande à chaque Métier}
\label{sec:c08-metiers}
% ============================================================================

\begin{boxmetier}{Analyste en réponse à incident}
Ne jamais suivre l'attaquant hors des systèmes du client~; demander par
écrit la conservation des données aux tiers dès le premier jour~; préparer
ses pièces comme si elles devaient être remises à un OPJ, ce qui arrive
souvent.
\end{boxmetier}

\begin{boxmetier}{Application OPJ/Expert Judiciaire}
L'OPJ requiert les tiers sans délai, en invoquant le dépérissement des
preuves (art.~54). L'expert s'en tient aux questions de sa mission et signale
à l'autorité tout fait découvert hors de celle-ci.
\end{boxmetier}

\begin{boxmetier}{RSSI}
Mettre la conservation des journaux en conformité avec l'article~29~: dix
ans pour les données de connexion et de trafic. La mesure se chiffre en
stockage, et elle conditionne toute enquête future.
\end{boxmetier}

\separateur

\begin{boxresume}
\textbf{Ce qu'il faut retenir de ce chapitre~:}
\begin{itemize}
    \item Les normes ISO/IEC~27037 à~27043 décrivent l'état de l'art de
          l'investigation~; l'ISO/IEC~27041 exige des méthodes validées.
    \item La loi n°~2010/012 confie l'enquête aux OPJ et aux agents de
          l'ANTIC (art.~52), encadre la saisie de données (art.~53, 54), le
          déchiffrement (art.~55) et l'expertise (art.~56).
    \item Les faits de NYANGA relèvent notamment des articles~65 al.~2,
          69, 71 et~72~; la qualification appartient au parquet.
    \item L'exploitant doit conserver dix ans ses données de connexion et
          de trafic (art.~29 al.~1)~: NYANGA ne les gardait que 60~jours.
    \item L'enquêteur privé n'a aucun pouvoir de contrainte~; l'entraide
          internationale passe par les autorités judiciaires.
\end{itemize}
\end{boxresume}

\begin{boxinfo}
\textbf{Points de contrôle.} Vous avez compris ce chapitre si vous savez~:
\begin{itemize}
    \item citer l'article qui fonde la saisie par copie de données~;
    \item rattacher un fait établi à une infraction et nommer l'élément à
          prouver~;
    \item dire ce qu'un enquêteur privé peut et ne peut pas faire.
\end{itemize}
\end{boxinfo}

% ============================================================================
\section*{Exercices}
% ============================================================================

\begin{boxexo}[title={\ding{45}\quad Exercice~8.1. Quel texte~? \niveaubase}]
Indiquez l'article de la loi n°~2010/012 qui répond à chaque question~:
(a) qui conduit l'enquête pénale en cas d'infraction cybernétique~?
(b) peut-on saisir une copie des données plutôt que le serveur~? (c) comment
obtenir la version en clair d'une archive chiffrée~? (d) combien de temps
l'exploitant doit-il conserver ses données de trafic~? (e) comment désigner
un expert~?
\end{boxexo}

\begin{boxexo}[title={\ding{45}\quad Exercice~8.2. Ce que l'enquêteur privé ne fait pas \niveaumoyen}]
Mme~MBARGA reçoit trois demandes de son client. Dites, pour chacune, si elle
peut y répondre et pourquoi~: (a) se connecter au serveur 203.0.113.47 avec
les identifiants trouvés dans l'outil de l'attaquant~; (b) examiner le
portable personnel de M.~KAMGA, qu'il apporte de lui-même~; (c) obtenir de
l'opérateur de paiement mobile l'identité des titulaires des onze
portefeuilles.
\end{boxexo}

\begin{boxexo}[title={\ding{45}\quad Exercice~8.3. Qualifier \niveaumoyen}]
Pour chacun des faits suivants, proposez un article de la loi n°~2010/012
et nommez l'élément que l'enquête doit établir~: (a) la création du compte
\texttt{svc\_backup} par l'attaquant, s'il était établi~; (b) la balise de
soixante secondes~; (c) l'effacement de l'archive \texttt{cpt.7z}~;
(d) l'envoi de 3,2~Go vers l'extérieur.
\end{boxexo}

\begin{boxexo}[title={\ding{45}\quad Exercice~8.4. Les obligations de la victime \niveauexpert}]
Un client dont le dossier figure parmi les fichiers exfiltrés assigne NYANGA
CRÉDIT en réparation.
\begin{enumerate}[(a)]
    \item Quels manquements aux articles~26, 29 et~30 peut-il invoquer, et
          sur quelles pièces de l'enquête~?
    \item NYANGA répond qu'elle est victime d'une infraction. Montrez que
          les deux positions ne s'excluent pas.
    \item Rédigez la recommandation du rapport d'enquête qui traite la
          conservation des journaux, en citant le texte.
\end{enumerate}
\end{boxexo}

\bigskip

\begin{boxinfo}
\textbf{Pour aller plus loin.}
\begin{itemize}
    \item La loi n°~2010/012~\cite{loi2010012}, dans son texte officiel.
    \item Les normes ISO/IEC~27037~\cite{iso27037_2012},
          27041~\cite{iso27041_2015}, 27042~\cite{iso27042_2015},
          27043~\cite{iso27043_2015}.
    \item La Convention de Budapest~\cite{budapest2001} et la Convention de
          Malabo~\cite{ua2014malabo}.
\end{itemize}
\end{boxinfo}

\begin{boxinfo}
\textbf{Transition.} L'enquête a montré par où l'attaquant est entré~: une
faille connue de la passerelle. La partie~IV change de point de vue~:
qu'est-ce qu'une vulnérabilité, et comment la mesurer~?
Chapitre~\ref{chap:vulnerabilite}.
\end{boxinfo}
